\begin{afterword}
\summaryhead

The essay argues that admitting your flaws is worth nothing unless you also plan to reduce them. What matters is progress.

It begins with Orthodox Judaism, as the author describes it. A proverb says each generation is to the one before as donkeys are to men, because law received at Sinai can only be lost in transmission. Modern rabbis therefore cannot overrule ancient ones, even on biology. As a child, the author concluded that science, which gains knowledge, must eventually surpass the Torah.

The Japanese phrase ``tsuyoku naritai,'' ``I want to become stronger,'' names the attitude the author prefers. The Yom Kippur litany Ashamnu, in which everyone confesses every sin, stands for the attitude he rejects: confession without improvement. He compares it to rationalists who announce that everyone is biased. A joke about a rabbi mocks pride in humility. The essay ends by demanding that anyone who confesses flaws also say what they will do about them.

\responsehead

The essay builds a rationalist identity against a religious caricature. Its one idea, that confessing a fault is empty without the resolve to fix it, is presented as a break with the tradition it describes. It is that tradition's own rule. Maimonides defines repentance as abandoning the sin and resolving never to repeat it. He compares confession without that resolve to immersing in a ritual bath while holding a dead lizard. The Yom Kippur liturgy includes a confession that asks ``that I may sin no more.'' The essay says the litany lacks exactly this, and lets the litany stand for the whole rite.

The other religious material is handled the same way. The angels-and-donkeys proverb is a conditional line of self-deprecation in the Talmud, and the essay turns it into a law of decline. The worm in the apple is permitted, in the Talmud, because of how a verse is read, not because of a theory of spontaneous generation. The Talmud's own rule for choosing between its sages favors the later generation, from Abaye and Rava on. At each point the essay reads its source loosely and in its own favor. The flattering conclusion follows: science progresses, religion decays, and the reader is on the right side.

The structure shows the same looseness. The thesis arrives in the fourth paragraph as a proverb. The title phrase appears in the fifth and is not used again until the last word. The opponent in the ninth paragraph, who beats a fist on the heart and says ``we are all biased,'' is never named or quoted.

The essay's closing rule fails when applied to the essay. ``Never confess to me that you are just as flawed as I am unless you can tell me what you plan to do about it.'' The essay contains no plan to reduce any bias. Two years later, the author wrote that the systematic training such a plan would need still did not exist.

In short: a rule the tradition already states, presented as a rebuke to that tradition, supported by misread sources and an unnamed opponent, and broken by the essay itself.
\end{afterword}
